The preceding section demonstrated that the tensorised formulation reproduces the angular power spectra computed by \texttt{CCL} to sub-per-cent accuracy across all observables and bin combinations relevant for the LSST Y1 survey scenario. Having established this numerical fidelity, we now turn to the practical implications of the framework for Stage-IV cosmological inference. This section addresses four complementary aspects: the computational performance of the tensorised implementation relative to conventional Limber codes (Section~\ref{sec:6.1}), the prospect of end-to-end emulation of the cosmology-dependent coefficients (Section~\ref{sec:6.2}), the capacity for analytic marginalisation over redshift distribution uncertainties (Section~\ref{sec:6.3}), and natural extensions to non-flat geometries and beyond-Limber scales (Section~\ref{sec:6.4}).

\subsection{Computational performance} \label{sec:6.1}

A central motivation for the tensorised reformulation is the clean separation of cosmology-dependent computation from the survey-specific projection. To quantify the resulting performance characteristics, we benchmark four implementations of the full \(3\times2\,\mathrm{pt}\) angular power spectrum computation: (i) \texttt{CCL}, which evaluates the Limber integral by numerical quadrature; (ii) \texttt{Numba--CPU}, a \texttt{LimberCloud} implementation accelerated with \texttt{Numba} just-in-time (JIT) compilation on CPU; (iii) \texttt{JAX--CPU}, the \texttt{LimberCloud} implementation compiled and executed on a multi-core CPU via \texttt{JAX}; and (iv) \texttt{JAX--GPU}, the same \texttt{JAX} implementation executed on a single NVIDIA GPU. All four implementations adopt identical cosmological parameters, multipole grids, tomographic binning, and redshift distributions, ensuring a direct comparison.

For the three \texttt{LimberCloud} implementations, the computation decomposes naturally into three clearly separated stages. The \textit{cosmology stage} computes the comoving distance--redshift relation \(\chi(z)\) and the non-linear three-dimensional power spectrum \(P_{\delta\delta}(k, z)\) using \texttt{CAMB} via the \texttt{CCL} back-end. The \textit{coefficient stage} evaluates the interval tensor \(B_{uv}^{ij,n}(\ell)\) (Equation~\ref{eq:5.4}) for all tracer pairs, grid-point indices, and multipoles using the 21 analytic integral formulas, and sums the result into the basis tensor \(B_{uv}^{ij}(\ell)\) (Equation~\ref{eq:5.5}). The \textit{projection stage} contracts the basis tensor with the discretised redshift distributions via Equation~\ref{eq:5.6} to obtain the angular power spectra for all tomographic bin pairs. This three-stage decomposition is a direct consequence of the tensorised structure: the cosmology and coefficient stages depend only on the cosmological parameters and the redshift grid, while the projection stage depends only on the tomographic bin assignments \(\phi^a_i\) and \(\phi^b_j\). For \texttt{CCL}, no such separation exists; the line-of-sight integration and summation over tomographic bins are performed together within a single monolithic call.

The benchmark consists of a series of sequential evaluations, each with an independently drawn cosmology sampled uniformly within \(\pm 10\%\) of the fiducial parameter values listed in Section~\ref{sec:3.1}. Cumulative wall-clock times are recorded for each stage as the number of evaluations is varied from 100 to 1000 in ten linearly spaced steps. For the \texttt{JAX} implementations, all timings include the overhead of JIT compilation and, for the GPU configuration, host-to-device memory transfers; explicit device synchronisation is enforced after every computation to ensure that all GPU operations have completed before the timer is read.

Figure~\ref{fig:8} presents the benchmark results for the LSST Y1 survey configuration. The left column shows the shape--shape (cosmic shear) observable alone, while the right column shows the full \(3\times2\,\mathrm{pt}\) computation including all three tracer-pair combinations. Within each column, the four panels display, from top to bottom, the total cumulative time and its decomposition into the cosmology, coefficient, and projection stages. In the lower three panels, the \texttt{CCL} curve shows the total \texttt{CCL} time as a fixed reference, since the \texttt{CCL} pipeline does not decompose into these stages; this convention allows the cost of each \texttt{LimberCloud} stage to be compared directly against the full conventional computation. Both configurations exhibit qualitatively consistent scaling behaviour: the cosmology stage dominates the total budget, the coefficient and projection stages are substantially cheaper, and the GPU implementation delivers the largest gains in the projection step. The \(3\times2\,\mathrm{pt}\) case shows proportionally larger absolute times in the coefficient and projection stages, reflecting the increased number of tracer-pair and tomographic-bin combinations, but the relative performance ordering among implementations is preserved.

\begin{figure*}[!htbp]
    \centering
    \begin{subfigure}[t]{0.48\textwidth}
        \centering
        \includegraphics[width=\textwidth]{Figures/BENCHMARK_SINGLE_Y1.pdf}
        \caption{Shape--shape (cosmic shear) only.}
        \label{fig:8a}
    \end{subfigure}
    \hfill
    \begin{subfigure}[t]{0.48\textwidth}
        \centering
        \includegraphics[width=\textwidth]{Figures/BENCHMARK_TRIPLE_Y1.pdf}
        \caption{Full \(3\times2\,\mathrm{pt}\) computation.}
        \label{fig:8b}
    \end{subfigure}
    \caption{Cumulative wall-clock time as a function of the number of sequential evaluations for four implementations of the angular power spectrum computation in the LSST Y1 survey configuration. \textit{Left:} shape--shape (cosmic shear) observable only. \textit{Right:} full \(3\times2\,\mathrm{pt}\) computation including shape--shape, position--shape, and position--position tracer pairs. Within each column, the panels show, from top to bottom, the total time and its decomposition into the cosmology, coefficient, and projection stages. In each of the lower three panels, the \texttt{CCL} curve (circles) reproduces the total \texttt{CCL} time as a reference, since the conventional pipeline does not separate into these stages. All implementations use identical cosmological parameters, multipole grids, and redshift distributions; each evaluation draws an independent random cosmology within \(\pm 10\%\) of the fiducial values.}
    \label{fig:8}
\end{figure*}

The cosmology stage dominates the total computation time for all implementations, accounting for more than half of the wall-clock budget at 1000 evaluations. This stage is shared across all four implementations, as it relies on the same \texttt{CAMB} and \texttt{CCL} back-end; the tight clustering of the four curves in the cosmology panel confirms that all differences in total runtime arise from the non-cosmological stages. This observation has an important practical implication: any further acceleration of the total pipeline requires either replacing the cosmology computation with a fast surrogate---as discussed in Section~\ref{sec:6.2}---or adopting a fast--slow sampling strategy \citep{2013PhRvD..87j3529L, 2021JCAP...05..057T, 2015A&C....12...45Z} that avoids re-evaluating the cosmology stage when only nuisance parameters change.

The coefficient stage reveals the first substantial performance differences among implementations. The \texttt{JAX--GPU} implementation evaluates the 21 analytic integral formulas approximately five times faster than the \texttt{Numba--CPU} alternative at 1000 evaluations. This gain arises because the coefficient evaluation involves independent integrals over \(N_\mathrm{Z} = 350\) redshift intervals for each multipole and tracer pair, a workload that maps naturally onto the massively parallel execution model of GPU hardware.

The projection stage exhibits the most dramatic performance gains. For the full \(3\times2\,\mathrm{pt}\) configuration (Figure~\ref{fig:8b}), the \texttt{JAX--GPU} implementation completes the projection---contracting the basis tensor with \(\phi^{a}_{i}\) and \(\phi^{b}_{j}\) for all bin pairs, tracer pairs, and multipoles---in approximately \(3\,\mathrm{s}\) at 1000 evaluations, compared with approximately \(45\,\mathrm{s}\) for \texttt{JAX--CPU} and approximately \(700\,\mathrm{s}\) for \texttt{Numba--CPU}. The shape--shape configuration (Figure~\ref{fig:8a}) exhibits the same performance hierarchy at proportionally lower absolute cost, consistent with the reduced number of bin-pair combinations. In both cases, the projection step reduces to a batch of bilinear contractions (Equation~\ref{eq:5.6}), an operation for which GPU linear-algebra libraries are heavily optimised. The \texttt{JAX--GPU} projection is faster than the full \texttt{CCL} pipeline by more than three orders of magnitude and faster than \texttt{JAX--CPU} by approximately one order of magnitude.

In terms of total runtime at 1000 evaluations, the \texttt{JAX--GPU} implementation is approximately twice as fast as \texttt{CCL}, with the gain constrained by the shared cosmology bottleneck. However, this headline figure significantly understates the advantage of the tensorised structure in the setting most relevant for practical inference. In a typical Markov chain Monte Carlo (MCMC) analysis, the parameter space comprises both model parameters \(\boldsymbol{\vartheta} = (\boldsymbol{\Omega},\, \boldsymbol{\lambda})\)---encompassing cosmological and astrophysical parameters---and nuisance parameters \(\boldsymbol{\Xi}\) controlling the redshift distributions (Section~\ref{sec:6.3}). Modern samplers exploit the observation that nuisance-parameter proposals do not require re-evaluation of the cosmological computation: under the tensorised formulation, proposals that modify only the redshift distributions require only the inexpensive projection step (Equation~\ref{eq:5.6}), while the basis tensor \(B_{uv}^{ij}(\ell)\) is reused from the previous cosmological evaluation. In this fast-proposal regime, the effective cost per likelihood evaluation on \texttt{JAX--GPU} is reduced to the projection time alone---approximately \(3\,\mathrm{ms}\) per evaluation---representing a speed-up of more than three orders of magnitude over the full \texttt{CCL} pipeline. Even on CPU, the \texttt{JAX--CPU} projection achieves a speed-up of approximately two orders of magnitude. This separation of time-scales provides a natural and efficient interface for fast--slow sampling strategies, substantially reducing the overall cost of high-dimensional inference.

\subsection{Prospects for end-to-end emulation} \label{sec:6.2}

The benchmark results above establish that the cosmology stage---the evaluation of the comoving distance--redshift relation and the non-linear three-dimensional power spectrum via \texttt{CAMB} and \texttt{HMCode2020}---dominates the total wall-clock time for all implementations. Emulators for the three-dimensional matter power spectrum are already mature: tools such as \texttt{CosmoPower} \citep{2022MNRAS.511.1771S}, \texttt{BACCO} \citep{2021MNRAS.507.5869A}, \texttt{EuclidEmulator2} \citep{2021MNRAS.505.2840E}, and the \texttt{Mira--Titan} emulator \citep{2017ApJ...847...50L} predict \(P_{\delta\delta}(k,z)\) to sub-per-cent accuracy in milliseconds, effectively eliminating the Boltzmann-solver bottleneck. A complementary factorisation strategy is the \texttt{COBRA} framework of \citet{2025PhRvL.134s1002B}, which decomposes the three-dimensional matter power spectrum itself into a small set of cosmology-independent basis functions whose amplitudes depend only on cosmological parameters; this approach operates at the level of \(P_{\delta\delta}(k,z)\) and is naturally combinable with the present framework, which factorises the projected angular power spectrum one step further along the survey axis and could in principle adopt a \texttt{COBRA}-style surrogate in place of the \texttt{CAMB} back-end in the cosmology stage. The tensorised structure of \texttt{LimberCloud} extends this acceleration one step further: rather than emulating only the three-dimensional power spectrum, it enables \textit{end-to-end emulation} of the basis tensor \(B_{uv}^{ij}(\ell)\), bridging the gap between three-dimensional emulation and the projected summary statistics that enter the likelihood.

The basis tensor depends on the full set of model parameters \(\boldsymbol{\vartheta} = (\boldsymbol{\Omega},\, \boldsymbol{\lambda})\), comprising both cosmological parameters \(\boldsymbol{\Omega}\) (e.g.\ \(\Omega_\mathrm{m}\), \(\sigma_8\), \(w_0\), \(w_a\)) and astrophysical parameters \(\boldsymbol{\lambda}\) (e.g.\ intrinsic-alignment amplitude, galaxy bias values, baryonic feedback strength) that enter through the three-dimensional power spectra \(P_{uv}(k,z)\). The emulation strategy proceeds in two phases of offline precomputation. In the first phase, an existing surrogate provides \(P_{\delta\delta}(k,z;\,\boldsymbol{\Omega})\) at negligible cost, eliminating the Boltzmann-solver and halo-model evaluations that dominate the cosmology stage. In the second phase, the emulated power spectrum is combined with the astrophysical bias models to form \(P_{uv}(k,z;\,\boldsymbol{\vartheta})\), and the coefficient stage evaluates \(B_{uv}^{ij}(\ell;\,\boldsymbol{\vartheta})\) via the 21 analytic integral formulas for a large training ensemble of parameter vectors \(\boldsymbol{\vartheta}\). A lightweight machine-learning surrogate is then trained to predict \(B_{uv}^{ij}(\ell)\) directly for any new \(\boldsymbol{\vartheta}\). At inference time, the emulator forward pass replaces both the cosmology and coefficient stages, and the full angular power spectrum for arbitrary tomographic bin pairs is reconstructed via the bilinear contraction of Equation~\ref{eq:5.6}.

The feasibility of this precomputation can be assessed from the benchmark results of Section~\ref{sec:6.1}. With the first-phase surrogate in place, only the coefficient stage needs to be evaluated directly, at a cost of approximately \(0.5\,\mathrm{s}\) per parameter vector on a single GPU for the full \(3\times2\,\mathrm{pt}\) configuration. A training set of \(10^6\) parameter vectors---of the same order as the number of independent basis-tensor elements---therefore requires approximately six days on a single GPU node. The workload is embarrassingly parallel across independent parameter vectors: distributing over \(\mathcal{O}(10)\) single-GPU nodes reduces the wall-clock time to approximately \(14\) hours.

In the most demanding scenario identified in Section~\ref{sec:5}, where scale- and bin-dependent biases require a separate basis tensor for each tomographic bin pair, the coefficient cost increases by up to a factor of \(\sim 25\) for the affected tracer pairs. For \(10^6\) training vectors this amounts to approximately \(100\) days on a single GPU, but the cost is readily reduced by combining parallelisation across \(\mathcal{O}(10)\) nodes, utilisation of multiple GPUs per node (up to four on typical HPC hardware), and explicit multi-process parallelism within each node in addition to the internal multi-threaded execution already employed by \texttt{JAX}. Together, these strategies can bring the wall-clock time to the order of one day. While the simplest and most widely used bias parameterisations in current Stage-IV analyses fall within the first two regimes of Section~\ref{sec:5}---for which the basis tensor is computed once per tracer pair---analyses incorporating beyond-linear, scale-dependent galaxy bias models that vary across tomographic selections would enter this third regime. The estimates above confirm that even this worst-case scenario remains accessible to the emulation framework.

The primary advantage of end-to-end emulation is a dramatic reduction in the cost per likelihood evaluation during inference. The cosmology and coefficient stages, which together account for the vast majority of the wall-clock time (Figure~\ref{fig:8}), are replaced by a neural-network forward pass whose cost is independent of the complexity of the underlying physical model. The remaining projection step executes in approximately \(3\,\mathrm{ms}\) per evaluation on GPU hardware (Section~\ref{sec:6.1}). The combined emulation-plus-projection pipeline is therefore expected to achieve a speed-up of three to four orders of magnitude over the conventional \texttt{CCL} approach, enabling rapid exploration of high-dimensional parameter spaces in both likelihood-based and simulation-based inference.

A further advantage is scalability. The projection step and the emulator forward pass are both embarrassingly parallel across parameter samples and tomographic bin pairs, so the throughput scales linearly with available compute units---additional GPUs translate directly into proportionally more evaluations per unit time. By contrast, the conventional \texttt{CCL} pipeline performs the Limber integration monolithically for each bin pair at each parameter-space point, offering limited scope for hardware-level parallelism beyond internal thread-level concurrency.

The main challenge for end-to-end emulation lies not in the dimensionality of the input parameter space---which is typically \(10\)--\(50\) and is well sampled by \(10^6\) training vectors---but in the dimensionality of the emulation target: the basis tensor comprises \(\sim\!10^6\) independent elements, each of which must be predicted to sufficient accuracy across all multipoles and tracer pairs simultaneously. The tensorised representation offers several structural advantages that mitigate this challenge. A large fraction of basis-tensor elements are exactly zero due to the sparsity of the kernel coefficients: for the \(r_\phi \times r_\phi\) case, the basis tensor is tridiagonal in \((i,j)\) and more than \(99\%\) of its entries vanish; for the mixed \(r_\phi \times r_\kappa\) case, the lower-triangular half is zero; and for the \(r_\kappa \times r_\kappa\) case, the matrix is dense but symmetric, so only the upper triangle contains independent information. These sparsity patterns and symmetries can be exploited through dimensionality-reduction techniques---such as principal component analysis (PCA) applied to the non-zero entries across the training set---to compress the emulation target to a low-dimensional latent space before training the surrogate model. Additional strategies include hierarchical or probe-specific emulator architectures that exploit the block structure across tracer pairs, and physics-informed network designs that encode known scaling relations among the basis-tensor components. A full demonstration of end-to-end emulation within a differentiable inference pipeline is left to follow-up work.

\subsection{Analytic marginalisation over redshift distribution uncertainties} \label{sec:6.3}

Complementary to the emulation strategy of the preceding subsection---which accelerates the cosmology-dependent computation---the tensorised formulation also enables a fundamentally different treatment of the nuisance parameters controlling the survey-specific redshift distributions. In conventional analyses, uncertainties in the tomographic redshift distributions are explored by sampling: at each MCMC step, the distributions are perturbed and the Limber integrals are re-evaluated in full. The bilinear structure of Equation~\ref{eq:5.6} makes it possible to marginalise over these uncertainties \textit{analytically}, eliminating the need to sample over them and reducing the effective dimensionality of the parameter space.

To formulate the marginalisation, we recast Equation~\ref{eq:5.6} in compact matrix notation. For a given tracer pair \((u,v)\) and multipole \(\ell\), define the basis-tensor matrix \([\boldsymbol{B}_{uv}(\ell)]_{ij} \equiv B_{uv}^{ij}(\ell)\) and the distribution matrix \(\boldsymbol{\Phi} \equiv [\boldsymbol{\phi}^1,\,\boldsymbol{\phi}^2,\,\ldots,\,\boldsymbol{\phi}^{N_\mathrm{B}}]\), an \((N_\mathrm{Z}{+}1) \times N_\mathrm{B}\) matrix whose \(a\)-th column is the discretised redshift distribution for tomographic bin \(a\), with \(N_\mathrm{B}\) denoting the number of tomographic bins. The full set of angular power spectra for all bin pairs is then expressed as the matrix bilinear form
\begin{equation} \label{eq:6.1}
    \boldsymbol{C}_{uv}(\ell) = \boldsymbol{\Phi}^{\mathrm{T}}\, \boldsymbol{B}_{uv}(\ell)\, \boldsymbol{\Phi},
\end{equation}
where \([\boldsymbol{C}_{uv}(\ell)]_{ab} = C_{uv}^{ab}(\ell)\) is the full matrix of angular power spectra across tomographic bins. This expression is exact and equivalent to Equation~\ref{eq:5.6}; its bilinear structure in the distribution matrix is the key property that enables the analytic treatment developed below.

We model uncertainties in the tomographic redshift distributions as perturbations about a fiducial model. Writing \(\boldsymbol{\Phi} = \bar{\boldsymbol{\Phi}} + \boldsymbol{\Xi}\), where \(\bar{\boldsymbol{\Phi}}\) is the fiducial distribution matrix and \(\boldsymbol{\Xi} \equiv [\boldsymbol{\xi}^1,\, \boldsymbol{\xi}^2,\, \ldots,\, \boldsymbol{\xi}^{N_\mathrm{B}}]\) is the perturbation matrix whose \(a\)-th column \(\boldsymbol{\xi}^a\) represents the departure of the true distribution from the fiducial in bin \(a\), the nuisance parameters over which we seek to marginalise are precisely the perturbation components \(\xi^a_i\). Substituting into Equation~\ref{eq:6.1} and expanding algebraically yields
\begin{equation} \label{eq:6.2}
    \boldsymbol{C}_{uv}(\ell;\,\boldsymbol{\Xi}) = \bar{\boldsymbol{C}}_{uv}(\ell) + \boldsymbol{\Xi}^{\mathrm{T}}\, \boldsymbol{B}_{uv}(\ell)\, \bar{\boldsymbol{\Phi}} + \bar{\boldsymbol{\Phi}}^{\mathrm{T}}\, \boldsymbol{B}_{uv}(\ell)\, \boldsymbol{\Xi} + \boldsymbol{\Xi}^{\mathrm{T}}\, \boldsymbol{B}_{uv}(\ell)\, \boldsymbol{\Xi},
\end{equation}
where \(\bar{\boldsymbol{C}}_{uv}(\ell) \equiv \bar{\boldsymbol{\Phi}}^{\mathrm{T}}\, \boldsymbol{B}_{uv}(\ell)\, \bar{\boldsymbol{\Phi}}\) is the fiducial angular power spectrum matrix. This expansion is \textit{exact}: because the angular power spectrum is bilinear in the distribution vectors, the polynomial terminates at second order in the perturbations and involves no approximation. This stands in contrast to existing analytic marginalisation approaches based on Taylor expansions of the Limber integral \citep{2020JCAP...10..056H, 2023MNRAS.522.5037R, 2023OJAp....6E..23H, 2025A&A...693A.178R, 2025arXiv250600758B, 2026arXiv260424425R}, which truncate the dependence on the nuisance parameters at first order and thereby introduce a systematic truncation error.

Collecting all angular power spectra across tracer pairs, bin pairs, and multipoles into a theory data vector \(\boldsymbol{\mathcal{M}}\), the expansion (Equation~\ref{eq:6.2}) yields the data-vector decomposition
\begin{equation} \label{eq:6.3}
    \boldsymbol{\mathcal{M}}(\boldsymbol{\vartheta},\, \boldsymbol{\Xi}) = \bar{\boldsymbol{\mathcal{M}}}(\boldsymbol{\vartheta}) + \boldsymbol{\Gamma}(\boldsymbol{\vartheta})\, \boldsymbol{\Xi} + \boldsymbol{\Xi}^{\mathrm{T}} \boldsymbol{\mathcal{B}}(\boldsymbol{\vartheta})\, \boldsymbol{\Xi},
\end{equation}
where \(\boldsymbol{\vartheta}\) denotes the model parameters (cosmological and astrophysical), \(\bar{\boldsymbol{\mathcal{M}}}(\boldsymbol{\vartheta})\) is the fiducial data vector, \(\boldsymbol{\Gamma}(\boldsymbol{\vartheta})\) is the linear response matrix encoding the first-order dependence on \(\boldsymbol{\Xi}\), and \(\boldsymbol{\Xi}^{\mathrm{T}} \boldsymbol{\mathcal{B}}(\boldsymbol{\vartheta})\, \boldsymbol{\Xi}\) denotes the vector whose entries are the bilinear contributions \((\boldsymbol{\xi}^{a})^{\mathrm{T}}\, \boldsymbol{B}_{uv}(\ell;\,\boldsymbol{\vartheta})\, \boldsymbol{\xi}^b\) for each data-vector element. Both terms follow directly from the algebraic structure of the bilinear form and require no finite-difference approximations. In practice, the high-dimensional nuisance vector \(\boldsymbol{\Xi}\) may be projected onto a low-dimensional basis describing physically motivated modes such as coherent redshift shifts and broadening; the formalism applies without modification to any such linear reparameterisation.

We now consider the analytic marginalisation of \(\boldsymbol{\Xi}\) under a Gaussian likelihood for the observed data vector \(\boldsymbol{\mathcal{D}}\),
\begin{equation} \label{eq:6.4}
    -2\ln\mathcal{L}(\boldsymbol{\mathcal{D}} \,|\, \boldsymbol{\vartheta},\, \boldsymbol{\Xi}) = \bigl[\boldsymbol{\mathcal{D}} - \boldsymbol{\mathcal{M}}(\boldsymbol{\vartheta},\, \boldsymbol{\Xi})\bigr]^{\mathrm{T}} \boldsymbol{\Sigma}^{-1} \bigl[\boldsymbol{\mathcal{D}} - \boldsymbol{\mathcal{M}}(\boldsymbol{\vartheta},\, \boldsymbol{\Xi})\bigr],
\end{equation}
where \(\boldsymbol{\Sigma}\) is the data covariance matrix. Defining the fiducial residual \(\boldsymbol{\Delta}(\boldsymbol{\vartheta}) \equiv \boldsymbol{\mathcal{D}} - \bar{\boldsymbol{\mathcal{M}}}(\boldsymbol{\vartheta})\) and substituting Equation~\ref{eq:6.3}, the quadratic form in the exponent separates into a Gaussian part and higher-order corrections,
\begin{align} \label{eq:6.5}
    \bigl[\boldsymbol{\mathcal{D}} - \boldsymbol{\mathcal{M}}\bigr]^{\mathrm{T}} \boldsymbol{\Sigma}^{-1} \bigl[\boldsymbol{\mathcal{D}} - \boldsymbol{\mathcal{M}}\bigr]
    ={} & \underbrace{\boldsymbol{\Delta}^{\mathrm{T}}\, \boldsymbol{\Sigma}^{-1}\, \boldsymbol{\Delta}}_{\mathcal{O}(\boldsymbol{\Xi}^0)} - \underbrace{2\,\boldsymbol{\Delta}^{\mathrm{T}}\, \boldsymbol{\Sigma}^{-1}\, \boldsymbol{\Gamma}\boldsymbol{\Xi}}_{\mathcal{O}(\boldsymbol{\Xi}^1)} + \underbrace{\boldsymbol{\Xi}^{\mathrm{T}}\, \boldsymbol{\Gamma}^{\mathrm{T}}\, \boldsymbol{\Sigma}^{-1}\, \boldsymbol{\Gamma}\, \boldsymbol{\Xi}}_{\mathcal{O}(\boldsymbol{\Xi}^2)} \nonumber\\[4pt]
    & \underbrace{-\, 2\,\boldsymbol{\Delta}^{\mathrm{T}}\, \boldsymbol{\Sigma}^{-1}\, (\boldsymbol{\Xi}^{\mathrm{T}} \boldsymbol{\mathcal{B}}\, \boldsymbol{\Xi})}_{\mathcal{O}(\boldsymbol{\Xi}^2)} + \underbrace{2\,\boldsymbol{\Xi}^{\mathrm{T}}\, \boldsymbol{\Gamma}^{\mathrm{T}}\, \boldsymbol{\Sigma}^{-1}\, (\boldsymbol{\Xi}^{\mathrm{T}} \boldsymbol{\mathcal{B}}\, \boldsymbol{\Xi})}_{\mathcal{O}(\boldsymbol{\Xi}^3)} \nonumber\\[4pt]
    & + \underbrace{(\boldsymbol{\Xi}^{\mathrm{T}} \boldsymbol{\mathcal{B}}\, \boldsymbol{\Xi})^{\mathrm{T}}\, \boldsymbol{\Sigma}^{-1}\, (\boldsymbol{\Xi}^{\mathrm{T}} \boldsymbol{\mathcal{B}}\, \boldsymbol{\Xi})}_{\mathcal{O}(\boldsymbol{\Xi}^4)}.
\end{align}
The first line collects terms of zeroth through second order in \(\boldsymbol{\Xi}\) and constitutes a Gaussian quadratic form in the nuisance vector; the subsequent lines arise from cross-terms and self-interactions involving the bilinear contribution \(\boldsymbol{\Xi}^{\mathrm{T}} \boldsymbol{\mathcal{B}}\, \boldsymbol{\Xi}\), spanning second through fourth order.

At linear order in the data-vector perturbation---neglecting the bilinear term \(\boldsymbol{\Xi}^{\mathrm{T}} \boldsymbol{\mathcal{B}}\, \boldsymbol{\Xi}\)---the integrand of the marginalisation integral is exactly Gaussian in \(\boldsymbol{\Xi}\). Assuming a Gaussian prior \(p(\boldsymbol{\Xi}) = \mathcal{N}(\boldsymbol{0},\, \boldsymbol{\Pi})\), the marginalised likelihood
\begin{equation} \label{eq:6.6}
    \mathcal{L}(\boldsymbol{\mathcal{D}} \,|\, \boldsymbol{\vartheta}) = \int \mathcal{L}(\boldsymbol{\mathcal{D}} \,|\, \boldsymbol{\vartheta},\, \boldsymbol{\Xi})\, p(\boldsymbol{\Xi})\, \mathrm{d}\boldsymbol{\Xi}
\end{equation}
is evaluated by completing the square in \(\boldsymbol{\Xi}\) and performing the Gaussian integral, yielding
\begin{equation} \label{eq:6.7}
    -2\ln\mathcal{L}(\boldsymbol{\mathcal{D}} \,|\, \boldsymbol{\vartheta}) = \boldsymbol{\Delta}^{\mathrm{T}}\, \tilde{\boldsymbol{\Sigma}}^{-1}(\boldsymbol{\vartheta})\, \boldsymbol{\Delta} + \ln\det\tilde{\boldsymbol{\Sigma}}(\boldsymbol{\vartheta}),
\end{equation}
where the effective covariance is
\begin{equation} \label{eq:6.8}
    \tilde{\boldsymbol{\Sigma}}(\boldsymbol{\vartheta}) = \boldsymbol{\Sigma} + \boldsymbol{\Gamma}(\boldsymbol{\vartheta})\, \boldsymbol{\Pi}\, \boldsymbol{\Gamma}(\boldsymbol{\vartheta})^{\mathrm{T}}.
\end{equation}
Equations~\ref{eq:6.7} and~\ref{eq:6.8} constitute the \textit{linear-order} analytic marginalisation result: the marginalisation over redshift distribution uncertainties reduces to an inflation of the data covariance by \(\boldsymbol{\Gamma}\boldsymbol{\Pi}\boldsymbol{\Gamma}^{\mathrm{T}}\), where \(\boldsymbol{\Gamma}\) is obtained directly from the basis tensor and the fiducial distributions. This result is structurally identical to that of existing first-order approaches \citep{2020JCAP...10..056H, 2023MNRAS.522.5037R, 2023OJAp....6E..23H, 2025A&A...693A.178R, 2025arXiv250600758B, 2026arXiv260424425R}; the key distinction is that \(\boldsymbol{\Gamma}\) is read off from the exact bilinear expansion (Equation~\ref{eq:6.2}) rather than estimated by finite differencing of the Limber integral.

We note that the bilinear contribution \(\boldsymbol{\Xi}^{\mathrm{T}} \boldsymbol{\mathcal{B}}\, \boldsymbol{\Xi}\), neglected in the linear-order result above, introduces non-Gaussian terms in the marginalisation integral through the second line of Equation~\ref{eq:6.5}. While these higher-order corrections can in principle be evaluated systematically---for instance using Isserlis' theorem \citep{1918Bka....12..134I} to compute the required polynomial moments of the Gaussian prior---they are expected to be small when the prior covariance \(\boldsymbol{\Pi}\) is well constrained by photometric redshift calibration, as the corrections are suppressed by additional powers of \(\lVert \boldsymbol{\xi}^a \rVert / \lVert \bar{\boldsymbol{\phi}}^a \rVert\) relative to the linear-order result. The tensorised formulation nevertheless provides these higher-order terms from first principles, offering a controlled framework in which the adequacy of the linear-order approximation can be assessed quantitatively. A detailed investigation of these contributions, together with numerical validation within a full inference pipeline, is left to follow-up work.

Compared to existing analytic marginalisation approaches, the tensorised formulation offers three key advantages. First, the dependence of the data vector on the nuisance parameters is known \textit{exactly} from the bilinear structure of the Limber projection---it is not an approximation obtained by Taylor-expanding the Limber integral. The complete polynomial structure, including the quadratic term, is accessible from first principles for systematic assessment.

Second, the linear response matrix \(\boldsymbol{\Gamma}(\boldsymbol{\vartheta})\) retains the full dependence on the cosmological and astrophysical parameters \(\boldsymbol{\vartheta}\), inherited from the cosmology dependence of the basis tensor \(\boldsymbol{B}_{uv}(\ell;\,\boldsymbol{\vartheta})\). In existing treatments, the effective covariance \(\tilde{\boldsymbol{\Sigma}}\) is typically evaluated at a single fiducial cosmology and held fixed throughout the inference, an approximation that may suppress degeneracies between redshift distribution uncertainties and cosmological or astrophysical parameters. The tensorised approach, particularly when combined with the end-to-end emulation of Section~\ref{sec:6.2}, preserves the cosmology dependence of \(\boldsymbol{\Gamma}\) at every point in parameter space, providing a principled baseline against which the fixed-cosmology approximation can be validated. In practice, this effect may be subdominant when the prior on the nuisance parameters is well constrained by external calibration data; the framework nevertheless provides the means to verify this assumption quantitatively.

Third, the entire marginalisation procedure operates on the precomputed basis tensor and does not require re-evaluation of the Limber projection, making it fully compatible with both emulated and directly evaluated pipelines. When combined with end-to-end emulation, analytic marginalisation further reduces the effective dimensionality of the sampling problem, enabling the MCMC chain to explore only the model parameters \(\boldsymbol{\vartheta}\) while the redshift distribution nuisance parameters are integrated out at the likelihood level.

\subsection{Future extensions} \label{sec:6.4}

The preceding subsections demonstrated that the tensorised factorisation enables substantial computational acceleration (Section~\ref{sec:6.1}), a natural pathway to end-to-end emulation (Section~\ref{sec:6.2}), and exact analytic marginalisation over redshift distribution uncertainties (Section~\ref{sec:6.3}). These capabilities rely on the bilinear structure of the Limber projection (Equation~\ref{eq:5.6}) and the clean separation of cosmology-dependent coefficients from survey-specific distributions. We now outline three directions in which the framework can be extended, each of which preserves this factorisation and therefore inherits the same emulation and marginalisation infrastructure.

The most direct generalisation concerns non-flat geometries. Throughout this work, we have adopted the spatially flat assumption \(f_K(\chi) = \chi\), which simplifies the lensing efficiency kernel and the multipole--wavenumber mapping. Although current CMB and baryon acoustic oscillation constraints restrict \(|\Omega_K|\) to the per-mille level \citep{2020A&A...641A...6P}, a self-consistent treatment of spatial curvature is desirable for analyses that explore extended parameter spaces or marginalise over \(\Omega_K\). For non-zero curvature, the transverse comoving distance \(f_K(\chi)\) (Equation~\ref{eq:2.2}) replaces \(\chi\) in the denominator of the Limber integral (Equation~\ref{eq:2.1}) and in the lensing kernel (Equation~\ref{eq:2.4}). The piecewise-linear strategy extends naturally to this setting: within each redshift interval \([\chi_n, \chi_{n+1}]\), the ratio \(f_K(\chi)/\chi\) varies slowly and can be linearised alongside the existing kernel and power-spectrum representations. The resulting integrands retain the rational-plus-logarithmic structure of the flat-sky case, with additional terms whose closed-form evaluation follows by the same methods used in the Appendix. Crucially, the linearised \(f_K(\chi)\) depends only on cosmology and the redshift grid, so the factorisation into a cosmology-dependent basis tensor and survey-specific distributions is fully preserved, and only a modest increase in the number of analytic integral formulas per interval is expected.

The kernel decomposition of Section~\ref{sec:4} is readily extensible to additional projected observables that share the same line-of-sight integral structure, as noted in Section~\ref{sec:2}. CMB lensing convergence provides the most immediate example: its radial kernel is structurally identical to \(W_\kappa^a(\chi)\) (Equation~\ref{eq:2.4}) with the source distribution replaced by a delta function at the last-scattering surface, so it is already accommodated by the existing \(r_\kappa\) coefficients without any modification to the integral formulas. Cross-correlating CMB lensing with photometric galaxy clustering and cosmic shear extends the standard \(3\times2\,\mathrm{pt}\) framework to a \(6\times2\,\mathrm{pt}\) analysis, an approach increasingly adopted in current surveys \citep{2025A&A...702A.169S, 2025A&A...699A.127J}; the tensorised formulation provides a natural and efficient infrastructure for such joint analyses. The integrated Sachs--Wolfe (ISW) effect introduces a radial kernel proportional to the time derivative of the gravitational potential, which can likewise be piecewise-linearised on the same redshift grid, adding a new set of basis-tensor components without modifying the existing machinery. Redshift-space distortions, which involve velocity-field kernels with a different radial weighting and angular-dependent terms, represent a further avenue for extension; their treatment requires additional care due to the derivative structure of the kernel and the anisotropy of the projection, but the underlying piecewise strategy remains applicable. Each of these extensions adds new kernel types to the existing framework while preserving the bilinear factorisation and its associated computational advantages.

A more challenging extension is the incorporation of exact angular power spectra beyond the Limber approximation. The Limber approximation is known to introduce per-cent-level errors at low multipoles (\(\ell \lesssim 20\)--\(50\)), particularly for galaxy clustering and galaxy--galaxy lensing where the radial kernels have narrow support \citep{2017A&A...602A..72C, 2025A&A...699A.124R}. The full (non-Limber) projection involves spherical Bessel functions \(j_\ell(k\chi)\) whose rapid oscillations preclude the simple rational-function integrands that underpin the present framework. Two complementary strategies offer promising routes forward. The first builds directly on the piecewise formalism: within each redshift interval, the kernel is piecewise linear and the power spectrum is piecewise linear in \(\chi\), so the integrand takes the form of a low-order polynomial multiplied by a spherical Bessel function over a finite interval. Integrals of this type admit closed-form solutions in terms of generalised hypergeometric functions and spherical Bessel functions evaluated at the interval endpoints \citep{2018PhRvD..97b3504G}; the same strategy can therefore be applied interval by interval, with the sparsity structure of the kernel coefficients exploited as in the Limber case. The second strategy employs FFTLog decompositions \citep{2000MNRAS.312..257H} to expand the power spectrum as a sum of power laws, \(P(k) = \sum_m c_m\, k^{\nu_m}\). When combined with the piecewise-linear kernel representation, each term in the expansion yields integrals involving products of power laws and Bessel functions that are known analytically \citep{2017JCAP...11..054A, 2020JCAP...05..010F, 2023PhRvD.108j3007F, 2025JOSS...10.8618R}. This hybrid approach offers improved numerical stability at low multipoles where the Bessel oscillations are most problematic, while retaining the separability of the tensorised framework. In either case, the key factorisation property---that the cosmology-dependent coefficients and the survey-specific distributions enter as separate factors in a bilinear contraction---would be preserved, enabling the same emulation (Section~\ref{sec:6.2}) and analytic marginalisation (Section~\ref{sec:6.3}) strategies developed above. When combined with the non-flat geometry and additional-probe extensions, a beyond-Limber generalisation of \texttt{LimberCloud} would provide a unified, tensorised forward model for joint multi-probe analyses spanning photometric and spectroscopic surveys together with CMB observations, offering a scalable foundation for next-generation cosmological inference.
